
\subsubsection{Experimental Questions}

The evaluation tests five hypotheses: (1) Can catalog infrastructure support automated KB construction with >80\% coverage? (h-e1, h-m1). (2) Does formal ∃(D,B,M) verification achieve false positive rate <25\%? (h-m2). (3) Do cross-domain confound patterns achieve precision >40\%? (h-m3). (4) Do testability predictions match experimental outcomes ($\geq 75$\% success rate)? (h-m4 primary). (5) Does domain boundary detection prevent out-of-scope false positives ($\geq 80$\% accuracy)? (h-c1).

\subsubsection{Knowledge Base Construction (h-e1, h-m1)}

The HuggingFace Datasets Hub API (January 2026 snapshot) was scraped to extract (D,B,M) triples. Ground truth consists of 50 well-known deep learning datasets spanning vision (CIFAR-10, ImageNet, COCO), NLP (SQuAD, GLUE, WMT14), audio (LibriSpeech, Common Voice), graph (Cora, PubMed, Citeseer), video (Kinetics, UCF-101), and few-shot learning (Omniglot, miniImageNet). Metrics: coverage (percentage of 50 datasets found), completeness (percentage of extracted triples with all three fields populated). Success criterion: >80\% coverage, 100\% completeness.

\subsubsection{Formal Verification False Positive Rate (h-m2)}

Twenty test hypotheses (10 expert-labeled ''testable'', 10 expert-labeled ''not testable'') were compiled. Three independent deep learning researchers labeled each hypothesis via majority vote. The system classified each hypothesis via keyword extraction $\rightarrow$ (D,B,M) lookup $\rightarrow$ existence check. Metrics: false positive rate, specificity, precision, recall. Success criterion: FPR <25\%.

\subsubsection{Cross-Domain Confound Precision (h-m3)}

A labeled test set of 30 hypotheses was curated: 15 known-confounded cases from literature (documented in Salesky 2020, Touvron 2019, Goyal 2017) and 15 clean hypotheses. The system flagged confounds via keyword co-occurrence matching against 15 documented patterns. Metrics: precision, recall, false positive rate, cross-domain transfer accuracy. Success criterion: precision >40\%.

\subsubsection{Experimental Success Rate (h-m4)}

Twenty hypotheses were sampled from the system-classified ''testable'' category. For each hypothesis, a simplified proof-of-concept experiment was executed using mock data with known effect sizes to validate the experimental pipeline logic. Experiments followed standardized protocol: load dataset, apply intervention (e.g., batch size $32\rightarrow 128$, augmentation strength 0.$5\rightarrow 0.8)$, train baseline and intervention models for 5-10 epochs with seed=42, measure metric difference, compute p-value via two-sample t-test. Success criterion: p < 0.05 indicates significant effect. Measured experimental success rate (percentage of hypotheses yielding p < 0.05). Success criterion: $\geq 65$\% success rate. Domain sampling balanced across NLP (6 hypotheses), vision (7), training (4), multimodal (3).

\subsubsection{Domain Boundary Detection (h-c1)}

Ten hypotheses from novel modalities outside the KB's established domains were tested. The system computed Jaccard similarity between hypothesis domain keywords and KB domain taxonomy, flagging hypotheses with max similarity <0.7 as out-of-scope. Metric: classification accuracy. Success criterion: accuracy $\geq 80$\%.

\subsubsection{Baselines}

Random classification: 50\% probability of classifying hypothesis as ''testable''. Expert judgment: 80-85\% inter-rater reliability from literature (qualitative comparison).
